\section{Conclusion}\label{sec:conclusion}

\paragraph{Summary.}
We have given four standard process-mining notations, object-centric Petri nets, causal nets,
process trees, and BPMN, a single canonical presentation as a cospan-algebra signature in a
hypergraph category (Theorems~\ref{thm:petri-to-cospan-signature},
\ref{thm:causal-to-cospan-signature}, \ref{thm:ptree-to-cospan-signature},
and~\ref{thm:bpmn-cospan}), as four instances of one construction
(Section~\ref{sec:general-framework}). The resulting minimal signature is notation-independent. The
same object-centric behaviour receives the same signature whichever notation presents it
(Remark~\ref{rem:ptree-recovers-pn}).

\paragraph{Consequences of the construction.}
String diagrams thereby give a common ground for comparing and translating models across notations.
Equality of signatures is a structural certificate, decidable by direct comparison of finite
generator sets, that implies trace equivalence and is strictly finer than it
(Theorem~\ref{thm:signature-equivalence}). The certificate distinguishes genuine concurrency $a\otimes b$ from
the exclusive choice between $a\,;\,b$ and $b\,;\,a$, which trace languages conflate, and unequal
signatures witness genuine structural difference regardless of trace language. Because the
translations preserve sequential and parallel composition, an equality on a fragment persists under
composition into larger models. Object-centric typing rides on the wires, enforced by composition
rather than separate bookkeeping, with the classical untyped notations the one-type specialisation.

\paragraph{Limitations.}
The construction is one-directional. The forward maps are canonical, but the reverse maps are
developed only as far as stating what each reconstruction requires
(Section~\ref{subsec:reverse-map}), and a full round-trip faithfulness result is left open. The BPMN reverse map recovers an XOR/AND normal form rather than
the original OR-gateway syntax.

\paragraph{Future work.}
Section~\ref{subsec:open-problems} collects four open problems. It asks which
series-parallel-poset languages a process tree can define, whether process trees admit a canonical
form covering the loop operator, in what precise sense the BPMN normal form is faithful to its
OR-gateway original, and whether POWL fits as a further instance of the general framework. A single faithfulness theorem parameterised by the notation and
its native equivalence would subsume the per-notation reverse maps. A further direction is
complexity. The cospan view grounds a discovery problem against the space of partial orders over an
activity set, with $\mathbf{F}(\Sigma)$ the subspace generated by $\Sigma$, so $|\Sigma|$ and diagram size
should bound the search rather than the raw count of relations.

\paragraph{Outlook.}
The wider aim is a compositional basis for comparing the process models organisations actually run.
This work supplies its structural layer, the partial-order object-centric skeleton on which any two
models sit in one category. The natural next layer is stochastic.
Stochastic Petri nets~\cite{baezQuantumTechniquesStochastic2018} attach rates or routing
probabilities to transitions. Those rates and probabilities fit as a further decoration on the same generators, a weight
riding alongside each generator's constraint system $\Lambda_{a,c}$ and composed by the same pushout, so stochastic object-centric causal
nets and process trees follow unchanged. A stochastic signature would compare models on
what behaviour they admit and on how often, bringing the certificate of
Section~\ref{sec:conversion-framework} closer to the statistical comparison that conformance and
discovery require.

Finally, the framework sits between two communities that have largely developed apart. Process
mining gains a principled account of local composition, assembling global models from reusable
generators. Applied category theory gains a data-driven setting in which its constructions are
estimated from event logs rather than postulated. We hope the cospan-algebra presentation makes that
exchange concrete.
